% =====================================================================
%  20.4  Lazy on Flattened Subtrees                    sec:c15-trees
%  Exports: prop:c15-flatten  fig:c15-subtree
%           lst:c15-620-soln  prob:c15-nytree
%  Reuses lst:c15-rupdate / lst:c15-rquery / lst:c15-assign (tag = paint).
%  Ch.5 owns tin/tout (thm:c24-range, def:c24-flatten, sec:c24-flatten) --
%  cited by label, NEVER re-derived.  Shared tree T (opener): subtree(2) is the
%  single Euler interval [2,5] holding nodes {2,4,5,8}.
%  CF 620E (live-verified on codeforces.com): tree n,m<=4e5, root 1, colours
%  1..60; type 1 v c = recolour subtree(v) to c; type 2 v = # distinct colours
%  in subtree(v).  Node record = 60-bit colour MASK, merge = OR, assign sets one
%  bit; query = popcount of the OR over [tin(v), tout(v)].
%  Sample 1 (n=7,m=10) -> 2 3 4 5 1 2  (hand-replayed).
%  CF 383C / CF 877E are Ch.18's difference-BIT subtree tasks, re-treated here
%  IN PROSE only (never re-linked) as the lazy-segtree twins that discharge the
%  non-invertible hand-off.
% =====================================================================
\section{Lazy on Flattened Subtrees}\label{sec:c15-trees}

\begin{previouslybox}[Ch.~5 flattened the tree onto a line; Ch.~20 finally updates that line]
Chapter~5 (\chref{c24}) proved the reduction this whole section stands on -- preorder times \(\tin/\tout\) turn ``the subtree of \(v\)'' into one contiguous interval \([\tin_v,\tout_v]\) of a flat array (\(\thmref{c24-range}\), \(\defref{c24-flatten}\)).  Chapter~18 rode that flat array with a Fenwick tree for the \emph{invertible} subtree tasks; Chapter~19 (\(\secref{c13-static}\)) noted the segment-tree twin and \emph{deferred the dynamic engine to here}.  This section supplies that engine: a lazy tree over the Euler array, so you can \emph{paint} or \emph{add to a whole subtree} in \(\bigO{\log n}\).  Nothing about \(\tin/\tout\) is re-derived -- only the range \emph{update} on the flattened line is new.
\end{previouslybox}

\begin{topicintro}
Given a rooted tree with a value at each node, ``add \(x\) to every node in \(v\)'s subtree'' or ``recolour \(v\)'s subtree'' or ``how many distinct colours are in \(v\)'s subtree?'' all become \emph{one range operation} on the array \(E\) defined by \(E[\tin_v]=\) (value at \(v\)): the subtree is \(E[\tin_v..\tout_v]\).  Build a lazy segment tree over \(E\).  The pending tag is exactly whatever \(\secref{c15-core}\)--\(\secref{c15-assign}\) taught (add, or paint); the only genuinely new ingredient is the \emph{node record}.  For a distinct-colour count, a node cannot store a scalar -- it stores the \term{set} of colours in its interval, as a 60-bit \term{bitmask}, merged by \(\mathbf{OR}\) (associative, identity \(0\)), and an assign collapses the mask to a single bit.  The answer is \(\mathit{popcount}\) of the \(\mathbf{OR}\) over the subtree's range.

\textbf{Prerequisites.} Ch.~5's \(\thmref{c24-range}\) (\(\tin/\tout\), subtree\(=\)interval), used by label.  \(\secref{c15-assign}\)'s paint tag and \(\secref{c15-core}\)'s \code{push}/\code{pull}.  The bitmask/merge idea is the Ch.~19 \emph{tuple aggregate} (\(\secref{c13-combine}\)) with \(\mathbf{OR}\) in place of \(+\).

\textbf{After this section you can}
\begin{itemize}[nosep]
  \item turn a subtree update/query into one range op on the Euler array by citing \(\thmref{c24-range}\); \hfill\lv{3}--\lv{4}
  \item store a colour \emph{set} as a bitmask at each node, OR-merge it, and collapse it under assign; \hfill\lv{4}
  \item explain Ch.~18's 383C / 877E as the same flattened tree seen through a lazy tree instead of a BIT; \hfill\lv{4}
  \item know that \emph{path} updates are the HLD boundary (Ch.~22) and are only decomposed, not coded, here. \hfill\lv{4}--\lv{5}
\end{itemize}
\end{topicintro}

% ---------------------------------------------------------------
\subsection{The reduction, cited not re-proved}
% ---------------------------------------------------------------

\begin{proposition}{A subtree lazy op is one range lazy op on the Euler array}{c15-flatten}
Let \(\tin/\tout\) be Ch.~5's preorder times on a rooted tree and \(E[\tin_v]\) the value at \(v\).  Then (i) ``update every node of \(\subtree(v)\) by a lazy-able transform \(g\)'' equals ``\(\code{rupdate}\) the interval \([\tin_v,\tout_v]\) by \(g\)'' on the lazy tree over \(E\); and (ii) ``aggregate of \(\subtree(v)\)'' equals ``\(\code{rquery}\) over \([\tin_v,\tout_v]\)''.  Both cost \(\bigO{\log n}\).
\end{proposition}
\begin{proofbox}
By \(\thmref{c24-range}\) the vertices of \(\subtree(v)\) are exactly those with \(\tin\) in the closed interval \([\tin_v,\tout_v]\), and the mapping vertex \(\mapsto\tin\) is a bijection onto that interval's indices.  Therefore applying a range operation to indices \([\tin_v,\tout_v]\) touches each subtree vertex once and no other vertex -- identical to the subtree operation, and the lazy tree's \(\bigO{\log n}\) bound (\(\thmref{c15-lazy}\)) applies unchanged.  The aggregate case is the query half of the same bijection.  Nothing about computing \(\tin/\tout\) is done here -- that is Ch.~5's DFS.
\end{proofbox}

\begin{figure}[H]
\centering
\begin{tikzpicture}[diagram, scale=0.62, transform shape, x=13mm, y=11mm]
  \node[boxnode, successgreen, minimum size=7mm, fill=successgreen!16] (1) at (0,4) {\(1\)};
  \node[boxnode, primaryblue, minimum size=7mm] (2) at (-1.7,3) {\(2\)};
  \node[boxnode, primaryblue, minimum size=7mm] (3) at (1.7,3) {\(3\)};
  \node[boxnode, successgreen, minimum size=7mm, fill=successgreen!16] (4) at (-2.6,2) {\(4\)};
  \node[boxnode, successgreen, minimum size=7mm, fill=successgreen!16] (5) at (-0.9,2) {\(5\)};
  \node[boxnode, primaryblue, minimum size=7mm] (6) at (1.0,2) {\(6\)};
  \node[boxnode, primaryblue, minimum size=7mm] (7) at (2.4,2) {\(7\)};
  \node[boxnode, successgreen, minimum size=7mm, fill=successgreen!16] (8) at (-0.9,1) {\(8\)};
  \node[boxnode, primaryblue, minimum size=7mm] (9) at (2.4,1) {\(9\)};
  \node[boxnode, primaryblue, minimum size=7mm] (10) at (2.4,0) {\(10\)};
  \draw[flow] (1)--(2); \draw[flow] (1)--(3);
  \draw[flow] (2)--(4); \draw[flow] (2)--(5);
  \draw[flow] (3)--(6); \draw[flow] (3)--(7);
  \draw[flow] (5)--(8); \draw[flow] (7)--(9); \draw[flow] (9)--(10);
  \node[annot, font=\scriptsize, anchor=west] at (3.2,2.5) {\(\subtree(2)=\{2,4,5,8\}\) (green)};
  % Euler strip E[1..10]
  \foreach \p/\nd in {1/1,2/2,3/4,4/5,5/8,6/3,7/6,8/7,9/9,10/10} {
    \draw[draw=gray!60, fill=gray!8] (\p-4.5,-1.2) rectangle (\p-3.5,-0.4);
    \node[font=\scriptsize\ttfamily] at (\p-4,-0.8) {\(\nd\)};
    \node[font=\tiny\ttfamily, text=darkgray] at (\p-4,-1.45) {\(\p\)};}
  \draw[difforange, line width=1pt] (-2.5,-0.25) -- (-2.5,-1.55);
  \draw[difforange, line width=1pt] (1.5,-0.25) -- (1.5,-1.55);
  \draw[difforange, line width=1pt] (-2.5,-0.25) -- (1.5,-0.25);
  \node[font=\scriptsize, text=difforange, anchor=south] at (-0.5,-0.2) {\([\tin_2,\tout_2]=[2,5]\)};
\end{tikzpicture}
\caption{The shared tree \(T\) (opener, Ch.~12--17).  Preorder gives \(\tin=(1,2,6,3,4,7,8,5,9,10)\), so \(\subtree(2)=\{2,4,5,8\}\) is exactly the \emph{contiguous} Euler slice \(E[2..5]=(2,4,5,8)\) -- the orange bracket.  Painting \(2\)'s subtree is one lazy range update on those four cells; no recursion into the tree, only Ch.~5's precomputed \(\tin/\tout\) and \(\coderef{c15-assign}\)'s paint.}
\label{fig:c15-subtree}
\end{figure}

% ---------------------------------------------------------------
\subsection{The distinct-colour record: a bitmask, OR-merged}
% ---------------------------------------------------------------

Colours are \(1..60\), so a set of colours fits a 64-bit word: bit \(k\) set means colour \(k{+}1\) is present.  The three lazy ingredients become:

\begin{itemize}[nosep,leftmargin=1.4em]
  \item \textbf{leaf} \(E[i]\) (a single vertex of colour \(c\)): mask \(=1\!\ll\!(c-1)\).
  \item \textbf{merge} \(\mathcal{M}=\mathbf{OR}\): parent mask = left \(\mid\) right -- the union of the two children's colour sets.  Identity \(=0\) (empty set), so an out-of-range query branch contributes nothing.
  \item \textbf{assign} tag: painting a covered node to colour \(c\) sets its mask to the single bit \(1\!\ll\!(c-1)\) and stashes \(c\) -- \(\propref{c15-override}\) verbatim (a paint erases the old set).
  \item \textbf{answer}: \(\mathit{popcount}\) of the OR of masks over \([\tin_v,\tout_v]\).
\end{itemize}

\noindent This is \(\secref{c13-combine}\)'s ``node stores a record, merge in \(O(1)\)'' pattern, and it is why a subtree-colour problem needs a \emph{segment} tree, not a BIT: a BIT accumulates a single number (\(+\)), it cannot OR a set under an \emph{assignment}.  \(\mathbf{OR}\) is associative and idempotent but has no inverse, so the difference trick is also out -- squarely lazy territory.

\begin{notebox}[The Ch.~18 hand-off, discharged: 383C and 877E]
Chapter~18 solved two subtree tasks on this same flattened array with a \emph{difference} Fenwick tree, and closed by admitting they only worked because the update was invertible:
\begin{itemize}[nosep,leftmargin=1.4em]
  \item \textbf{CF 383C (Propagating tree)} -- ``add \(v\) to a node's whole subtree, but \emph{alternating sign} by depth parity, point read.''  Ch.~18 split it into two parity classes and used two BITs (an invertible \(\pm\) difference).  The lazy-segtree twin is the one most people reach for: one range-add on \([\tin_v,\tout_v]\), with the sign baked in by keeping \emph{two} sums per node (even- and odd-depth) and letting a push flip which child gets \(+d\) vs \(-d\).  Same answer, but the structure now also survives a \emph{non-invertible} variant for free.
  \item \textbf{CF 877E (Danil and a Part-time Job)} -- ``flip the lights in a subtree (a \emph{toggle}), query the sum of a subtree.''  Ch.~18 handled the toggle because a flip is its own inverse.  As a lazy tree this is textbook: node record \(=\) count of lit bulbs, tag \(=\) flip (\(\secref{c15-compose}\)), and a flip folds as \(\mathit{count}\leftarrow\mathit{len}-\mathit{count}\) with \(\mathit{flip}^2=\mathit{id}\) merging two flips.
\end{itemize}
Both are \emph{re-treated here in prose only} -- the URLs are Chapter~18's and are never re-linked (the anti-repetition rule).  What this section adds is that once your update stops being a plain add/flip (say, \emph{paint to colour \(c\)}, as in CF~620E), the BIT routes are closed and the lazy tree is the only game.
\end{notebox}

\paragraph{Path updates are Ch.~22's.}
``Update every node on the path \(u\!\to\!v\)'' is \emph{not} one subtree interval; heavy-light decomposition cuts the path into \(O(\log n)\) contiguous Euler segments, each of which is one lazy range op.  The chains themselves, and why \(O(\log n)\) of them, are Chapter~22 (\chref{c16}) -- forward-cited, with \emph{no} decomposition code here.  Subtree updates (this section) need no HLD; they are already a single interval.

% ---------------------------------------------------------------
\subsection{Problem Lab: CF 620E New Year Tree}
% ---------------------------------------------------------------

\solvedproblem{c15-nytree}{New Year Tree}{CF 620E}
  {https://codeforces.com/problemset/problem/620/E}{CF 2100}
  {Euler-flatten subtree; node = 60-bit colour mask, OR merge, assign collapses it}

\begin{problemstatement}[CF 620E New Year Tree]
A tree of \(n\) vertices rooted at \(1\); vertex \(i\) has colour \(c_i\in[1,60]\).  Process \(m\) queries: \texttt{1 v c} --- change every vertex in the subtree of \(v\) to colour \(c\); \texttt{2 v} --- print the number of \emph{distinct} colours present in the subtree of \(v\).

\textbf{Input.} \(1\le n,m\le 4\cdot10^{5}\); \(1\le c_i\le 60\); then \(n-1\) edges; then \(m\) queries as above.

\textbf{Output.} For each type-2 query, one line: the distinct-colour count.

\textbf{Example (official, first sample).}
\begin{verbatim}
Input:
7 10
1 1 1 1 1 1 1
1 2
1 3
1 4
3 5
3 6
3 7
1 3 2
2 1
1 4 3
2 1
1 2 5
2 1
1 6 4
2 1
2 2
2 3

Output:
2
3
4
5
1
2
\end{verbatim}
\end{problemstatement}

\begin{attemptgate}[Attempt it yourself first -- \lv{4}, give yourself 30 minutes]
``Distinct colours in a subtree'' is a set, not a number -- a BIT (Ch.~18) has nothing to accumulate.  And \(n\le 4\cdot10^5\) subtree repaints must not touch each vertex.
\peek{can a set of colours from \(1..60\) fit in one machine word?}
\nudge{\(60\le 64\), so a colour set is a 64-bit mask and union is \(\mathbf{OR}\).  Now the subtree is a range (\(\thmref{c24-range}\)); keep a lazy tree whose node value is that mask.}
\fullhint{Flatten with one DFS (\(\tin/\tout\)).  Segment tree over \(E[1..n]\): leaf mask \(=1\!\ll\!(c-1)\); internal node = left \(\mid\) right; assign tag = colour \(c\) (fold sets mask \(=1\!\ll\!(c-1)\), override the past, \(\propref{c15-override}\)).  Type~1 = \code{paint}\((\tin_v,\tout_v,c)\).  Type~2 = \(\mathit{popcount}\) of \code{query}\((\tin_v,\tout_v)\).  Colours \(\le60\) \(\Rightarrow\) \code{unsigned long long}; \(\mathit{OR}\)'s identity \(0\) for empty branches.}
\end{attemptgate}

\solstep{Understanding}
Two reductions stack.  Ch.~5 turns each subtree query into a range query on \(E\).  The distinct-colour requirement turns the scalar aggregate into a \emph{set}, but a \(1..60\) set is a 64-bit word, so the merge is bit-parallel \(\mathbf{OR}\) and the assignment collapses a node to one bit.  What remains is a range-assign / range-query lazy tree -- \(\secref{c15-assign}\) with the mask record swapped in.

\solstep{Key Observations}
\begin{observation}
\(\mathbf{OR}\) over colour masks is associative, commutative, idempotent, identity \(0\) -- a perfectly good segment-tree merge (\(\secref{c13-combine}\)) even though it is not a sum; a BIT cannot host it because there is no inverse and updates overwrite.
\end{observation}
\begin{observation}
A subtree is \emph{one} contiguous Euler interval, so no HLD is needed -- that machinery (\(\chref{c16}\)) is only for \emph{paths}.
\end{observation}
\begin{observation}
The paint tag is a colour, not a mask, so a covered node's mask becomes the single bit for that colour; keeping the mask as \(\mathit{node}\)-value and the colour as \(\mathit{tag}\) is enough -- you never store a mask in the tag.
\end{observation}

\solstep{Think It Through With Me}
\begin{thinkbox}
\deadend{Recolour the subtree by walking its vertices}
{A subtree can be \(\Theta(n)\) wide; \(m\le 4\cdot10^5\) such walks is \(\bigO{nm}\) -- \tle.  Flatten and paint as one range.}
\deadend{Store the colour \emph{count} per node}
{Two colours can have equal counts; a count cannot answer ``how many \emph{distinct} colours.''  Store the \emph{set} (a mask), union by \(\mathbf{OR}\), answer by \(\mathit{popcount}\).}
\aha{Flatten (Ch.~5) so subtree\(=\)range; let each node hold a 60-bit colour mask OR-merged; paint is a range-assign that collapses the mask to one bit.}
\end{thinkbox}

\solstep{The Approach Ladder}
\begin{center}\renewcommand{\arraystretch}{1.25}
\begin{tabular}{@{}c p{0.19\linewidth} p{0.26\linewidth} l c@{}}
\toprule
\textbf{\#} & \textbf{Approach} & \textbf{Idea} & \textbf{Time} & \textbf{Verdict}\\
\midrule
1 & \lv{3}\ walk the subtree & recolour/count by DFS each query & \(\bigO{nm}\) & \tle \\
2 & \lv{3}\ BIT of counts & cannot union distinct colours & -- & wrong \\
3 & \lv{4}\ \textbf{Euler + mask tree} & paint/OR/popcount on \(E\) & \(\bigO{(n+m)\log n}\) & \ac \\
4 & \lv{4}\ set-of-colours per node & \(\mathbf{OR}\) as \texttt{std::set} union & \(\bigO{(n{+}m)\log^2}\) & \ac, heavier \\
\bottomrule
\end{tabular}\end{center}
Rung~4 (a \texttt{std::set} per node) also works since colours \(\le60\), but the bitmask does union in one instruction -- the point of choosing a machine-word record.

\solstep{Solution}
One DFS fills \(\tin, \tout,\) and \(E\).  Build a lazy tree over \(E\) with leaf mask \(1\!\ll\!(c-1)\) and \(\mathbf{OR}\) merge.  Type~1: \code{paint} applied to \(\tin_v,\tout_v\) with colour \(c\) (an assign tag).  Type~2: print \(\mathit{popcount}\) of \code{query} on \(\tin_v,\tout_v\).

\solstep{Proof of Correctness}
\begin{proofbox}
\textbf{Flatten.} By \(\thmref{c24-range}\) the vertices of \(\subtree(v)\) have \(\tin\in[\tin_v,\tout_v]\) exactly, so operating on that index interval is operating on the subtree (\(\propref{c15-flatten}\)).
\textbf{Mask invariant.} Each node's mask has bit \(k\) iff some vertex in its Euler range currently has colour \(k{+}1\).  \emph{Build:} leaf is its vertex's single colour; internal \(=\) OR of children (union).  \emph{Paint:} a covered node set to colour \(c\) is uniformly \(c\), so its mask becomes the single bit \(1\!\ll\!(c-1)\) and the assign tag makes children agree when pushed (\(\propref{c15-override}\)); a partial paint \code{push}es then \code{pull}s, restoring the union from the children.
\textbf{Answer.} \code{query} ORs the masks of the canonical cover of \([\tin_v,\tout_v]\), a disjoint partition of the subtree's indices; by associativity/idempotence of \(\mathbf{OR}\) the union equals the set of colours in the subtree, whose size is \(\mathit{popcount}\).
\end{proofbox}

\solstep{Implementation}
\begin{lstlisting}[caption={CF 620E full solution: Euler-flatten + 60-bit colour mask, OR merge, assign tag}, label={lst:c15-620-soln}]
#include <bits/stdc++.h>
using namespace std;
typedef unsigned long long ull;
const int MX = 400005;
int n, m, col[MX], tin[MX], tout[MX], euler[MX], tmr;
vector<int> adj[MX];
ull t[4*MX]; int tagv[4*MX];          // tagv: colour owed, 0 = none

void dfs(int v, int p){
    tin[v] = ++tmr; euler[tmr] = v;
    for (int to : adj[v]) if (to != p) dfs(to, v);
    tout[v] = tmr;
}
void apply(int v, int c){
    t[v] = 1ULL << (c - 1);           // whole node painted to one colour
    tagv[v] = c;                       // assign overrides
}
void push(int v){
    if (!tagv[v]) return;
    apply(2*v, tagv[v]); apply(2*v+1, tagv[v]); tagv[v] = 0;
}
void build(int v, int l, int r){
    tagv[v] = 0;
    if (l == r){ t[v] = 1ULL << (col[euler[l]] - 1); return; }
    int m = (l + r) / 2; build(2*v,l,m); build(2*v+1,m+1,r);
    t[v] = t[2*v] | t[2*v+1];         // OR = union of colours
}
void paint(int v, int l, int r, int ql, int qr, int c){
    if (ql > r || qr < l) return;
    if (ql <= l && r <= qr){ apply(v, c); return; }   // cover: fold + stop
    push(v); int m = (l + r) / 2;
    paint(2*v,l,m,ql,qr,c); paint(2*v+1,m+1,r,ql,qr,c);
    t[v] = t[2*v] | t[2*v+1];                          // pull
}
ull query(int v, int l, int r, int ql, int qr){
    if (ql > r || qr < l) return 0;                     // OR identity = empty set
    if (ql <= l && r <= qr) return t[v];
    push(v); int m = (l + r) / 2;
    return query(2*v,l,m,ql,qr) | query(2*v+1,m+1,r,ql,qr);
}
int main(){
    scanf("%d%d", &n, &m);
    for (int i = 1; i <= n; i++) scanf("%d", &col[i]);
    for (int i = 0; i < n-1; i++){ int x,y; scanf("%d%d", &x,&y); adj[x].push_back(y); adj[y].push_back(x); }
    dfs(1, 0);
    build(1, 1, n);
    while (m--){
        int tk; scanf("%d", &tk);
        if (tk == 1){ int v,c; scanf("%d%d", &v,&c); paint(1,1,n, tin[v], tout[v], c); }
        else { int v; scanf("%d", &v);
               printf("%d\n", __builtin_popcountll( query(1,1,n, tin[v], tout[v]) )); }
    }
}
\end{lstlisting}

\solstep{Trace on the Official Sample}
\begin{tracebox}[CF 620E sample 1: preorder of \(1\!\to\!\{2,3,4\},\,3\!\to\!\{5,6,7\}\)]
DFS preorder \(E=(1,2,3,5,6,7,4)\) (all vertices start colour \(1\)); \(\tin=(1\!:\!1,2\!:\!2,3\!:\!3,5\!:\!4,6\!:\!5,7\!:\!6,4\!:\!7)\), \(\tout(3)=6\), \(\tout(1)=7\).\\
\textbf{1 3 2} paint \(\subtree(3)=[3,6]\) colour \(2\) \(\Rightarrow\) mask there \(=\) bit2. \textbf{2 1}: subtree \(1=[1,7]\) union \(\{\text{bit1,bit2}\}\) \(\Rightarrow\boxed{2}\).\\
\textbf{1 4 3} paint \([7,7]\) bit3. \textbf{2 1}: \(\{1,2,3\}\Rightarrow\boxed{3}\).\\
\textbf{1 2 5} paint \([2,2]\) bit5. \textbf{2 1}: \(\{1,2,3,5\}\Rightarrow\boxed{4}\).\\
\textbf{1 6 4} paint \([5,5]\) bit4. \textbf{2 1}: \(\{1,2,3,4,5\}\Rightarrow\boxed{5}\).\\
\textbf{2 2} \([2,2]=\) bit5 \(\Rightarrow\boxed{1}\). \textbf{2 3} \([3,6]=\{2,4\}\) (nodes \(3,5,6,7\) colours \(2,2,4,2\)) \(\Rightarrow\boxed{2}\).\\
Printed \(\boxed{2\;3\;4\;5\;1\;2}\) -- official.
\end{tracebox}

\solstep{Complexity}
\begin{complexitybox}
One DFS \(O(n)\); build \(O(n)\); each query \(\bigO{\log n}\) with \(O(1)\) 64-bit \(\mathbf{OR}\)/\code{apply}.  Total \(\bigO{(n+m)\log n}\approx 8\cdot10^5\cdot19\approx 1.5\cdot10^7\) node-visits for \(n,m\le4\cdot10^5\) -- \ac within 3 s.  Memory \(4n\) \code{ull} masks \(+\;4n\) \code{int} tags \(+\) adjacency \(\approx 30\) MB.
\end{complexitybox}

\solstep{Verification and Edge Cases}
\begin{center}\renewcommand{\arraystretch}{1.2}
\begin{tabular}{@{}p{0.22\linewidth} p{0.30\linewidth} p{0.36\linewidth}@{}}
\toprule
\textbf{Case} & \textbf{Why tricky} & \textbf{Code behaviour}\\
\midrule
colour \(60\) & bit shift overflow & \code{1ULL<<59}, \code{ull} mask (\(\ge64\) bits) \\
single-vertex subtree & leaf range & \(\tin_v{=}\tout_v\), one cell painted/read \\
re-paint same subtree & override & \code{apply} replaces mask and tag (\(\propref{c15-override}\)) \\
empty query branch & recursion clip & returns \(0\), \(\mathbf{OR}\) ignores it \\
root query \([1,n]\) & whole tree & root mask returned by cover-guard, \(O(1)\) \\
\bottomrule
\end{tabular}\end{center}

\begin{hardeningbox}
\begin{itemize}[nosep,leftmargin=1.4em]
  \item \textbf{Colours reach 60, so use \code{unsigned long long}.}  A 32-bit \code{int} mask truncates colours \(>32\); \(\mathit{popcount}\) then under-counts.  \code{1ULL << (c-1)}, never \code{1 << (c-1)}.
  \item \textbf{The DFS is recursive on up to \(4\cdot10^5\) nodes.}  A path-shaped tree overflows the default call stack -- either raise the stack limit or run the Euler tour iteratively (\chref{c24}'s own warning).  The segment-tree recursions are only \(\log n\) deep and are safe.
  \item \textbf{Paint, do not ``add a colour.''}  Type~1 \emph{sets} the subtree's colour, so the assign tag collapses the mask to one bit; if you OR-ed the new colour into the mask instead, stale colours never disappear and \(\mathit{popcount}\) only grows.
  \item \textbf{Read \([\tin_v,\tout_v]\), not \([\tin_v,\tout_v)\).}  Ch.~5's \(\tout\) here is the inclusive last preorder index of the subtree; mixing an exclusive \(\tout\) drops the deepest vertex.
\end{itemize}
\end{hardeningbox}

% ---------------------------------------------------------------
\subsection{Common Pitfalls}
% ---------------------------------------------------------------

\begin{warningbox}[Trap 1: re-deriving \(\tin/\tout\) here]
This section is \emph{allowed} to assume \(\thmref{c24-range}\).  A subtle bug is hand-rolling a slightly-off Euler tour (0- vs 1-based, inclusive vs exclusive \(\tout\)) when Ch.~5 already fixed the convention.  Cite it; compute the times once by the same DFS Ch.~5 prescribes.
\end{warningbox}

\begin{warningbox}[Trap 2: a scalar aggregate for a distinct-count question]
Storing ``number of vertices of colour \(c\)'' answers a \emph{count} query, not ``how many \emph{different} colours.''  Distinct-ness is a property of a \emph{set}; merge sets (\(\mathbf{OR}\)), don't add counts.
\end{warningbox}

\begin{warningbox}[Trap 3: forgetting to \code{push} before a partial subtree read]
A later \(\mathbf{OR}\) query over a range that only \emph{partially} covers a painted node reads that node's children -- which are stale until you push the paint down.  Same rule as \(\secref{c15-core}\); the fact that the aggregate is a mask does not exempt it.
\end{warningbox}

% ---------------------------------------------------------------
\subsection{Where to Drill Next}
% ---------------------------------------------------------------

\begin{stuckbox}
\begin{itemize}[nosep,leftmargin=1.4em]
  \item Did you build the Euler array \(E[\tin_v]=\) value at \(v\), and index the tree by \(\tin\), not by vertex id?
  \item Is your \(\mathit{OR}\) identity \(0\), and your mask at least 64 bits wide?
  \item Does a paint \emph{replace} the mask (single bit), not accumulate?
  \item Is \(\tout\) used inclusively, matching Ch.~5's definition?
  \item For a deep tree, is the DFS recursion (or its iterative substitute) safe?
\end{itemize}
\end{stuckbox}

Related in the mixed set: CF~1638E (\emph{Colorful Operations}) shifts colours along \emph{paths} -- you still flatten and range-update, but the ranges come from a different decomposition, so read it after Ch.~22; CF~755D (\emph{PolandBall and Polygon}) is a subtree-exclusion add/sum on a tour, the invertible cousin Ch.~18 could partly do.  The path version of \emph{this} problem (paint a root-to-node path lazily) is exactly Ch.~22's \(\chref{c16}\) HLD, where the same range-assign tag runs on each chain segment.

\paragraph{Back-cites and forwards.}
Chapter~5 (\chref{c24}) owns \(\tin/\tout\) and the flatten theorem -- cited, never re-proved (\(\secref{c24-flatten}\)).  Chapter~19 (\(\secref{c13-static}\)) owns the static aggregate/merge idea and the tuple-record pattern reused as the mask.  Chapter~18 (\chref{c14}) owns 383C/877E, re-treated above in prose.  Forward: path updates go to Ch.~22 (\chref{c16}); the non-composable updates stay out of the subtree engine entirely (Ch.~23, \(\secref{c15-strains}\)).

% ---------------------------------------------------------------
\begin{recallbox}[Section 20.4 in 10 seconds]
Ch.~5's \(\thmref{c24-range}\) makes \(\subtree(v)\) the Euler interval \([\tin_v,\tout_v]\); so a \emph{subtree} lazy update/query is a single \emph{range} op on the flattened array \(E\) (\(\propref{c15-flatten}\)).  For ``distinct colours'' the node record is a 60-bit \emph{mask}, \(\mathbf{OR}\)-merged (identity \(0\)), and assign collapses it to one bit; the answer is \(\mathit{popcount}\) over the range.  Path updates are Ch.~22's HLD, not this section.  Worked: CF~620E (\(\probref{c15-nytree}\), \(\coderef{c15-620-soln}\)); prose re-treats 383C/877E, discharging Ch.~18's hand-off.
\selftest{
  \item State, in one line, why a subtree recolour is a range assign and not a tree walk.
  \item Why must the mask be \(\ge64\) bits, and what breaks with 32?
  \item Which merge replaces ``\(+\)'' here, what is its identity, and why can't a BIT host it?
  \item Give the Ch.~18 task that this section re-treats in prose and name its lazy tag.
  \item Why does \emph{path} lazy need Ch.~22 while \emph{subtree} lazy does not?
}
\end{recallbox}
